\documentclass[11pt,a4paper]{article}

\usepackage[utf8]{inputenc}
\usepackage[T1]{fontenc}
\usepackage{geometry}
\usepackage{hyperref}
\usepackage{enumitem}
\usepackage{xcolor}
\usepackage{fancyhdr}
\usepackage{amsmath}
\usepackage{booktabs}
\usepackage{listings}

\geometry{margin=1in}
\hypersetup{
  colorlinks=true,
  linkcolor=blue,
  urlcolor=blue
}

\setlist{itemsep=0.35em, topsep=0.35em}

\pagestyle{fancy}
\fancyhf{}
\lhead{Object Oriented Programming (B.Tech CSE)}
\rhead{Unit 01: Introduction to OOPs}
\cfoot{\thepage}
\setlength{\headheight}{14pt}

\lstset{
  language=Java,
  basicstyle=\ttfamily\small,
  keywordstyle=\color{blue},
  commentstyle=\color{gray},
  stringstyle=\color{teal},
  showstringspaces=false,
  columns=fullflexible,
  frame=single
}

\begin{document}

\begin{center}
  {\LARGE \textbf{Assignment -- Unit 01}}\\[0.25em]
  {\large Object Oriented Programming (CSEG1044)}\\[0.25em]
  \normalsize (OOP principles; Java program structure; I/O; types; operators; control flow; arrays; strings)
\end{center}

\noindent
\textbf{Instructor:} Tofik Ali \hfill \textbf{Submission Deadline:} March 31, 2026\\
\textbf{Student Name:} \_\_\_\_\_\_\_\_\_\_\_\_\_\_\_\_\_ \hfill \textbf{Roll No.:} \_\_\_\_\_\_\_\_\\

\vspace{0.6em}
\hrule
\vspace{0.8em}

\textbf{Instructions}
\begin{itemize}[leftmargin=*]
  \item Answer \textbf{all} questions. Write assumptions where needed.
  \item There is \textbf{no time limit}. Submit on or before \textbf{March 31, 2026}.
  \item For programming questions, write \textbf{clean Java code}: meaningful variable names, small methods, and comments where needed.
  \item Unless a question says otherwise, use only the \textbf{standard Java library} (no external libraries).
  \item Each program should be in a separate \texttt{.java} file with the same name as the \texttt{public class} (if any).
  \item Where applicable, include a \textbf{sample input} and \textbf{sample output}.
\end{itemize}

\vspace{0.6em}

\section*{Easy (10 Questions)}
\begin{enumerate}[label=E\arabic*., leftmargin=*]
  \item \textbf{Key OOP terms (short):} Define (1--2 lines each): class, object, state, behavior, invariant.
    Give one example invariant for a \texttt{BankAccount} or \texttt{Student}.

  \item \textbf{OOP principles (match):} For each statement, identify the most relevant principle:
    \textbf{abstraction / encapsulation / inheritance / polymorphism}. (One principle per statement.)
    \begin{enumerate}[label=(\alph*), leftmargin=*]
      \item ``Hide the balance field and allow changes only through \texttt{deposit()} and \texttt{withdraw()}.''\\
      \item ``A \texttt{SavingAccount} is a specialized kind of \texttt{BankAccount}.''\\
      \item ``I only care that \texttt{print()} works; I don't care how it prints.''\\
      \item ``A variable of type \texttt{Shape} can refer to a \texttt{Circle} or \texttt{Rectangle}.''\\
    \end{enumerate}

  \item \textbf{Java ecosystem:} Explain JDK vs JRE vs JVM in 4--6 lines.
    Write one line: what is bytecode and why is it important?

  \item \textbf{Hello Java (first run):} Write a program \texttt{Hello.java} that prints:
    \texttt{Hello, UPES!}
    Then write the compile and run commands.

  \item \textbf{Command line arguments:} Write \texttt{ArgsEcho.java} that prints:
    \begin{itemize}[leftmargin=*]
      \item number of arguments,
      \item each argument with its index.
    \end{itemize}
    Example run: \texttt{java ArgsEcho A B C}.

  \item \textbf{Scanner input + formatting:} Write \texttt{SumAvg.java} that reads two numbers and prints sum and average using
    \texttt{printf} with 2 decimal places.

  \item \textbf{Choose correct data types:} For each, choose a suitable Java type and justify in one line:
    \begin{enumerate}[label=(\alph*), leftmargin=*]
      \item Student age
      \item CPI/CGPA (0--10, may contain decimals)
      \item A yes/no value
      \item A large population count (may exceed 2 billion)
    \end{enumerate}

  \item \textbf{Operator pitfalls (predict output):} Predict output (and write 1-line reason for each):
    \begin{enumerate}[label=(\alph*), leftmargin=*]
      \item \texttt{System.out.println(5/2);}
      \item \texttt{System.out.println(5/2.0);}
      \item \texttt{System.out.println(2 + 3 * 4);}
      \item \texttt{System.out.println((2 + 3) * 4);}
    \end{enumerate}

  \item \textbf{If--else (leap year):} Write \texttt{LeapYear.java} that reads a year and prints \texttt{LEAP} or \texttt{NOT LEAP}.
    (Use the correct leap year rules.)

  \item \textbf{Strings: \texttt{==} vs \texttt{equals}:} Write a small program that demonstrates why
    \texttt{==} is not used for string content comparison. Explain in 3--5 lines what happened.
\end{enumerate}

\section*{Medium (10 Questions)}
\begin{enumerate}[label=M\arabic*., leftmargin=*]
  \item \textbf{Array summary:} Write \texttt{MarksSummary.java}:
    read \texttt{n} marks into an array, then print min, max, sum, and average.

  \item \textbf{Second largest (distinct):} Write \texttt{SecondLargest.java} to find the second largest \textbf{distinct} element
    in an integer array. If it doesn't exist, print a clear message.

  \item \textbf{Palindrome (strings):} Write \texttt{Palindrome.java} that checks whether an input string is a palindrome.
    Ignore spaces and case. Example: \texttt{``Never odd or even''} $\rightarrow$ palindrome.

  \item \textbf{Vowel/consonant count:} Write \texttt{VowelCount.java} that reads a line and counts vowels, consonants,
    digits, spaces, and ``other'' characters.

  \item \textbf{Menu-driven calculator:} Write \texttt{CalculatorMenu.java}:
    display a menu (add/sub/mul/div/exit), repeat until exit, and handle divide-by-zero using \textbf{if} conditions.

  \item \textbf{Args-based conversion:} Write \texttt{TempConvert.java} that takes two args:
    a number and a mode (\texttt{C2F} or \texttt{F2C}). Print converted temperature with 2 decimals.
    Example: \texttt{java TempConvert 37 C2F}.

  \item \textbf{2D arrays (matrix addition):} Write \texttt{MatrixAdd.java} that reads two matrices of same size
    and prints their sum.

  \item \textbf{Remove duplicates (preserve order, no collections):} Given an integer array, create a new array
    that contains each value only once in the order of first appearance. Do not use \texttt{ArrayList/HashSet}.

  \item \textbf{StringBuilder join:} Write \texttt{JoinWords.java}:
    read \texttt{n} words and print them joined by comma (CSV style) using \texttt{StringBuilder}.
    Example output: \texttt{apple,banana,mango}.

  \item \textbf{Modelling exercise (design only):} Problem: ``A library stores books and members. Members can borrow/return books.''\\
    Write:
    \begin{itemize}[leftmargin=*]
      \item 4--6 candidate classes,
      \item 2 responsibilities per class,
      \item 2 invariants total (any classes),
      \item 2 short scenarios (use cases).
    \end{itemize}
\end{enumerate}

\section*{Hard (10 Questions)}
\begin{enumerate}[label=H\arabic*., leftmargin=*]
  \item \textbf{Student report manager (arrays):} Write \texttt{StudentReport.java}:
    store names and marks for \texttt{n} students using arrays.
    Provide menu options: add, update marks by name, search by name, show class average, show topper, exit.

  \item \textbf{Mini OOP class (intro):} Implement a \texttt{BankAccount} class with:
    account number, balance, and methods \texttt{deposit}, \texttt{withdraw}, \texttt{transferTo}.
    Enforce the invariant \texttt{balance >= 0} using validation.

  \item \textbf{Anagram checker:} Write \texttt{Anagram.java} that checks if two strings are anagrams.
    Ignore spaces and case. Example: \texttt{listen} and \texttt{silent}.

  \item \textbf{Character frequency (ASCII):} Write \texttt{CharFreq.java}:
    read a line and print frequency of each character that appears. (Hint: use an int array of size 256.)

  \item \textbf{Sorting (bubble sort):} Write \texttt{BubbleSort.java} to sort an integer array in ascending order
    without using built-in sort. Print the array before and after sorting.

  \item \textbf{Pattern printing:} Write a program to print this pattern for input \texttt{n}:
    \begin{lstlisting}[language=]
n=5
*
**
***
****
*****
    \end{lstlisting}

  \item \textbf{Guess the number (loops + conditions):} Write \texttt{GuessGame.java}:
    generate a random number between 1 and 100, keep asking user to guess, and print ``too high/too low'' until correct.
    Also print total attempts.

  \item \textbf{Date validation:} Write \texttt{ValidateDate.java} that reads \texttt{dd mm yyyy} and prints VALID/INVALID.
    Consider month lengths and leap years.

  \item \textbf{First non-repeating character:} Write \texttt{FirstUniqueChar.java}:
    given a string, find the first character that occurs exactly once.
    If none exists, print a message.

  \item \textbf{Explain your choices:} Pick any \textbf{two} programs from the Hard section and write 8--12 lines each:
    algorithm steps, edge cases, and why your approach works.
\end{enumerate}

\vfill
\noindent\textit{End of Assignment}

\end{document}

